\clearpage
\chapter{Conclusion}
\label{ch:conclusion}

This thesis asked two questions, how well bankruptcy can be predicted when the
models learn from real firms only and which financial features the fitted
models then rely on. The first question is answered by the framework itself.
With ClusterCentroids undersampling, repeated stratified cross-validation and
the $F_2$-score as the selection metric, gradient boosting reaches an $F_2$ of
$0.550$ and the random forest $0.546$, both clearly ahead of the logistic
regression at $0.482$ and above the comparable result reported by
\textcite{wangliu2021}.

The answer to the second question is that there is no single ranking. 
Although the three importance methods diverge, they seem to do so along 
their known weaknesses, impurity importance splits credit between redundant 
features, permutation hides features with close substitutes and SHAP shares 
the credit while at the same time keeping it visible. Among these differences 
one signal remains stable, as each model puts a measure of debt dependence 
either at the top or very close to it. The contribution of the thesis therefore 
lies in the comparison. If each ranking were considered only on its own it would 
already seem to give the correct answer, it is only by placing them side by side 
that one can see how much of the ranking is influenced by the model and the method 
behind it.

Both of these contributions indicate the same direction for future research. Since 
the framework is not tied to this dataset, the natural next step is to apply it to 
data from other countries and markets in order to find out whether the undersampling 
approach holds and whether similar ratios appear under the importance methods. 
Datasets whose collection is documented more transparently would additionally allow 
first steps toward causal statements, which this thesis avoided carefully. 
Such data is rare, as collecting firm records of this scope is complex and costly, 
which is exactly why the well-studied benchmark Taiwanese dataset was chosen here. 
Until then, the value of this thesis lies in a pipeline whose results can be trusted 
for what they are, predictions and associations measured on real firms.